% GENERATED FILE. Edit canonical lessons, then run npm run generate:books.
% canonical-source-hash: fnv1a64:66f8950527890590
% canonical-lessons: FR-C426-reception, FR-C426-receptionniste, FR-C426-etage, FR-C426-sejour, FR-C426-deballer

\chapter{À la réception}
\label{ch:fr-a-la-reception}
\hlchaptermodality{426}

\begin{chapteropening}
\textbf{By the end of this chapter:} \emph{I can check in at a hotel: the front desk, the receptionist, the floor, the stay and unpacking.}

Complete the last lesson of chapter 426: I can check in at a hotel: the front desk, the receptionist, the floor, the stay and unpacking.
\end{chapteropening}

\section[la réception]{la réception — the reception desk; the reception (party)}

\label{lesson:FR-C426-reception}



\begin{quote}
Before the new one: say the French for to camp, then the French for to bivouac.
\end{quote}

\subsection*{You'll want to know: la réception}
\textbf{la réception} — \textquotedblleft{}the reception desk; the reception (party)\textquotedblright{}.

\begin{quote}
Nous laissons la clé à la réception.
\end{quote}

\begin{quote}
\emph{We leave the key at reception.}
\end{quote}

\emph{Une réception} is also a formal party.

\begin{cousinweb}
\emph{réception} is a learned borrowing from Latin \emph{receptiō} 'a receiving', from \emph{recipere} 'to receive', built on \emph{capere} 'to take' (EtymDB).
\end{cousinweb}

\begin{grammarlens}[title={What you've built}]
One more word for telling, explaining and describing at B1.
\end{grammarlens}

\subsection*{Your turn}
\begin{itemize}
\raggedright
  \item \emph{Say it:} \emph{la réception}
  \item \emph{Say it:} \emph{la réception}, once more
  \item \emph{Say it:} a sentence of your own with \emph{la réception}
  \item \emph{Recall:} say the French for to camp, then the French for to bivouac, then say \emph{la réception} again
\end{itemize}

\begin{tcolorbox}[breakable,colback=teal!4,colframe=teal!35!black,title={Before you move on}]
What does la réception mean? (\textquotedblleft{}The reception desk; the reception (party)\textquotedblright{}.) Say it once more.
\end{tcolorbox}
\section[le réceptionniste]{le réceptionniste — the receptionist}

\label{lesson:FR-C426-receptionniste}



\begin{quote}
Before the new one: say the French for to bivouac, then the French for the reception desk.
\end{quote}

\subsection*{You'll want to know: le réceptionniste}
\textbf{le réceptionniste} — \textquotedblleft{}the receptionist\textquotedblright{}.

\begin{quote}
Le réceptionniste nous a donné un plan du quartier.
\end{quote}

\begin{quote}
\emph{The receptionist gave us a map of the area.}
\end{quote}

Same form with \emph{la} for a woman.

\begin{cousinweb}
\emph{réceptionniste} is built on \emph{réception} with \emph{-iste} (EtymDB lists it); \emph{réception} goes back to Latin \emph{receptiō}.
\end{cousinweb}

\begin{grammarlens}[title={What you've built}]
One more word for telling, explaining and describing at B1.
\end{grammarlens}

\subsection*{Your turn}
\begin{itemize}
\raggedright
  \item \emph{Say it:} \emph{le réceptionniste}
  \item \emph{Say it:} \emph{le réceptionniste}, once more
  \item \emph{Say it:} a sentence of your own with \emph{le réceptionniste}
  \item \emph{Recall:} say the French for to bivouac, then the French for the reception desk, then say \emph{le réceptionniste} again
\end{itemize}

\begin{tcolorbox}[breakable,colback=teal!4,colframe=teal!35!black,title={Before you move on}]
What does le réceptionniste mean? (\textquotedblleft{}The receptionist\textquotedblright{}.) Say it once more.
\end{tcolorbox}
\section[l'étage]{l'étage — the storey, the floor (of a building)}

\label{lesson:FR-C426-etage}



\begin{quote}
Before the new one: say the French for the reception desk, then the French for the receptionist.
\end{quote}

\subsection*{You'll want to know: l'étage}
\textbf{l'étage} — \textquotedblleft{}the storey, the floor (of a building)\textquotedblright{}.

\begin{quote}
Notre chambre est au troisième étage.
\end{quote}

\begin{quote}
\emph{Our room is on the third floor.}
\end{quote}

A masculine noun. \emph{Le rez-de-chaussée} is the ground floor; the \emph{premier étage} is one floor up.

\begin{cousinweb}
\emph{étage} comes from Old French \emph{estage} 'dwelling, level', which is usually traced to a Vulgar Latin word built on \emph{stāre} 'to stand'. EtymDB lists it without a source.
\end{cousinweb}

\begin{grammarlens}[title={What you've built}]
One more word for telling, explaining and describing at B1.
\end{grammarlens}

\subsection*{Your turn}
\begin{itemize}
\raggedright
  \item \emph{Say it:} \emph{l'étage}
  \item \emph{Say it:} \emph{l'étage}, once more
  \item \emph{Say it:} a sentence of your own with \emph{l'étage}
  \item \emph{Recall:} say the French for the reception desk, then the French for the receptionist, then say \emph{l'étage} again
\end{itemize}

\begin{tcolorbox}[breakable,colback=teal!4,colframe=teal!35!black,title={Before you move on}]
What does l'étage mean? (\textquotedblleft{}The storey, the floor (of a building)\textquotedblright{}.) Say it once more.
\end{tcolorbox}
\section[le séjour]{le séjour — the stay; the living room}

\label{lesson:FR-C426-sejour}



\begin{quote}
Before the new one: say the French for the receptionist, then the French for the storey.
\end{quote}

\subsection*{You'll want to know: le séjour}
\textbf{le séjour} — \textquotedblleft{}the stay; the living room\textquotedblright{}.

\begin{quote}
Le patron de l'hôtel nous a souhaité un bon séjour.
\end{quote}

\begin{quote}
\emph{The hotel owner wished us a pleasant stay.}
\end{quote}

\emph{La salle de séjour} or \emph{le séjour} — the living room.

\begin{cousinweb}
\emph{séjour} is built on \emph{séjourner} 'to stay', which goes back to Vulgar Latin \emph{subdiurnāre}, built on \emph{diēs} 'day' (Lewis \& Short).
\end{cousinweb}

\begin{grammarlens}[title={What you've built}]
One more word for telling, explaining and describing at B1.
\end{grammarlens}

\subsection*{Your turn}
\begin{itemize}
\raggedright
  \item \emph{Say it:} \emph{le séjour}
  \item \emph{Say it:} \emph{le séjour}, once more
  \item \emph{Say it:} a sentence of your own with \emph{le séjour}
  \item \emph{Recall:} say the French for the receptionist, then the French for the storey, then say \emph{le séjour} again
\end{itemize}

\begin{tcolorbox}[breakable,colback=teal!4,colframe=teal!35!black,title={Before you move on}]
What does le séjour mean? (\textquotedblleft{}The stay; the living room\textquotedblright{}.) Say it once more.
\end{tcolorbox}
\section[déballer]{déballer — to unpack}

\label{lesson:FR-C426-deballer}



\begin{quote}
Before the new one: say the French for the storey, then the French for the stay.
\end{quote}

\subsection*{You'll want to know: déballer}
\textbf{déballer} — \textquotedblleft{}to unpack\textquotedblright{}.

\begin{quote}
Il a déballé ses affaires et rangé ses chemises.
\end{quote}

\begin{quote}
\emph{He unpacked his things and put his shirts away.}
\end{quote}

Figurative: \emph{déballer tout} — to come out with everything.

\begin{cousinweb}
\emph{déballer} is built from \emph{dé-} and \emph{balle} 'bale, bundle', which is usually traced to a Frankish word meaning 'ball'. EtymDB does not list the verb.
\end{cousinweb}

\begin{grammarlens}[title={What you've built}]
One more word for telling, explaining and describing at B1.
\end{grammarlens}

\subsection*{Your turn}
\begin{itemize}
\raggedright
  \item \emph{Say it:} \emph{déballer}
  \item \emph{Say it:} \emph{déballer}, once more
  \item \emph{Say it:} a sentence of your own with \emph{déballer}
  \item \emph{Recall:} say the French for the storey, then the French for the stay, then say \emph{déballer} again
\end{itemize}

\begin{tcolorbox}[breakable,colback=teal!4,colframe=teal!35!black,title={Before you move on}]
What does déballer mean? (\textquotedblleft{}To unpack\textquotedblright{}.) Say it once more.
\end{tcolorbox}
